% GENERATED FILE. Edit canonical lessons, then run npm run generate:books.
% canonical-source-hash: fnv1a64:d7ad5a81d9606409
% canonical-lessons: TA-C438-kokku, TA-C438-kuyil, TA-C438-kiri, TA-C438-palli, TA-C438-tel, TA-R438-review-432-438

\chapter{Herons and Koels}
\label{ch:ta-herons-and-koels}
\hlchaptermodality{438}

\begin{chapteropening}
\textbf{By the end of this chapter:} \emph{I can talk about herons, koels, mongooses, lizards and scorpions.}

Complete the last lesson of chapter 438: I can talk about herons, koels, mongooses, lizards and scorpions.
\end{chapteropening}

\section[\texorpdfstring{\ta{கொக்கு} (kokku)}{கொக்கு (kokku)}]{\ta{கொக்கு} (kokku) — a heron; a crane}

\label{lesson:TA-C438-kokku}



\begin{quote}
Before the new one: say the Tamil for to hunt, then the Tamil for a species.
\end{quote}

\subsection*{You'll want to know: \ta{கொக்கு}}
\textbf{\ta{கொக்கு}} — \emph{kokku} — \textquotedblleft{}a heron; a crane\textquotedblright{}.

\begin{quote}
\ta{வயலில்} \ta{வெள்ளைக்} \ta{கொக்குகள்} \ta{நின்றன}.
\end{quote}

\begin{quote}
\emph{vayalil veḷḷaik kokkukaḷ niṉṟaṉa.}
\end{quote}

\begin{quote}
White herons stood in the field.
\end{quote}

The plural is \ta{கொக்குகள்}.

\begin{cousinweb}
\textbf{\ta{கொக்கு}} (\emph{kokku}) is native Dravidian (DEDR 2125), \textquotedblleft{}a crane, a stork, a paddy bird\textquotedblright{}.
\end{cousinweb}

\begin{grammarlens}[title={What you've built}]
One more word for getting around.
\end{grammarlens}

\subsection*{Your turn}
\begin{itemize}
\raggedright
  \item \emph{Say it:} \emph{kokku}
  \item \emph{Say it:} \emph{kokku}, once more
  \item \emph{Say it:} \emph{vayalil veḷḷaik kokkukaḷ niṉṟaṉa.} — White herons stood in the field.
  \item \emph{Recall:} say the Tamil for to hunt, then the Tamil for a species, then say \emph{kokku} again
\end{itemize}

\begin{tcolorbox}[breakable,colback=teal!4,colframe=teal!35!black,title={Before you move on}]
What does \ta{கொக்கு} mean? (\textquotedblleft{}A heron; a crane\textquotedblright{}.) Say it once more.
\end{tcolorbox}
\section[\texorpdfstring{\ta{குயில்} (kuyil)}{குயில் (kuyil)}]{\ta{குயில்} (kuyil) — a koel; an Indian cuckoo}

\label{lesson:TA-C438-kuyil}



\begin{quote}
Before the new one: say the Tamil for a species, then the Tamil for a heron.
\end{quote}

\subsection*{You'll want to know: \ta{குயில்}}
\textbf{\ta{குயில்}} — \emph{kuyil} — \textquotedblleft{}a koel; an Indian cuckoo\textquotedblright{}.

\begin{quote}
\ta{காலையில்} \ta{குயில்} \ta{இனிமையாகப்} \ta{பாடியது}.
\end{quote}

\begin{quote}
\emph{kālaiyil kuyil iṉimaiyākap pāṭiyatu.}
\end{quote}

\begin{quote}
In the morning the koel sang sweetly.
\end{quote}

Its call is heard most in the hot months.

\begin{cousinweb}
\textbf{\ta{குயில்}} (\emph{kuyil}) is native Dravidian (DEDR 1764), \textquotedblleft{}the koel, the Indian cuckoo\textquotedblright{}.
\end{cousinweb}

\begin{grammarlens}[title={What you've built}]
One more word for getting around.
\end{grammarlens}

\subsection*{Your turn}
\begin{itemize}
\raggedright
  \item \emph{Say it:} \emph{kuyil}
  \item \emph{Say it:} \emph{kuyil}, once more
  \item \emph{Say it:} \emph{kālaiyil kuyil iṉimaiyākap pāṭiyatu.} — In the morning the koel sang sweetly.
  \item \emph{Recall:} say the Tamil for a species, then the Tamil for a heron, then say \emph{kuyil} again
\end{itemize}

\begin{tcolorbox}[breakable,colback=teal!4,colframe=teal!35!black,title={Before you move on}]
What does \ta{குயில்} mean? (\textquotedblleft{}A koel; an Indian cuckoo\textquotedblright{}.) Say it once more.
\end{tcolorbox}
\section[\texorpdfstring{\ta{கீரி} (kīri)}{கீரி (kīri)}]{\ta{கீரி} (kīri) — a mongoose}

\label{lesson:TA-C438-kiri}



\begin{quote}
Before the new one: say the Tamil for a heron, then the Tamil for a koel.
\end{quote}

\subsection*{You'll want to know: \ta{கீரி}}
\textbf{\ta{கீரி}} — \emph{kīri} — \textquotedblleft{}a mongoose\textquotedblright{}.

\begin{quote}
\ta{கீரி} \ta{பாம்புடன்} \ta{சண்டை} \ta{போட்டது}.
\end{quote}

\begin{quote}
\emph{kīri pāmpuṭaṉ caṇṭai pōṭṭatu.}
\end{quote}

\begin{quote}
The mongoose fought with the snake.
\end{quote}

With -\ta{உடன்} for the opponent: \ta{பாம்புடன்} \ta{சண்டை} \ta{போடு}.

\begin{cousinweb}
\textbf{\ta{கீரி}} (\emph{kīri}) is native Dravidian (DEDR 1613), \textquotedblleft{}a mongoose\textquotedblright{}.
\end{cousinweb}

\begin{grammarlens}[title={What you've built}]
One more word for getting around.
\end{grammarlens}

\subsection*{Your turn}
\begin{itemize}
\raggedright
  \item \emph{Say it:} \emph{kīri}
  \item \emph{Say it:} \emph{kīri}, once more
  \item \emph{Say it:} \emph{kīri pāmpuṭaṉ caṇṭai pōṭṭatu.} — The mongoose fought with the snake.
  \item \emph{Recall:} say the Tamil for a heron, then the Tamil for a koel, then say \emph{kīri} again
\end{itemize}

\begin{tcolorbox}[breakable,colback=teal!4,colframe=teal!35!black,title={Before you move on}]
What does \ta{கீரி} mean? (\textquotedblleft{}A mongoose\textquotedblright{}.) Say it once more.
\end{tcolorbox}
\section[\texorpdfstring{\ta{பல்லி} (palli)}{பல்லி (palli)}]{\ta{பல்லி} (palli) — a lizard; a house gecko}

\label{lesson:TA-C438-palli}



\begin{quote}
Before the new one: say the Tamil for a koel, then the Tamil for a mongoose.
\end{quote}

\subsection*{You'll want to know: \ta{பல்லி}}
\textbf{\ta{பல்லி}} — \emph{palli} — \textquotedblleft{}a lizard; a house gecko\textquotedblright{}.

\begin{quote}
\ta{சுவரில்} \ta{ஒரு} \ta{பல்லி} \ta{ஓடியது}.
\end{quote}

\begin{quote}
\emph{cuvaril oru palli ōṭiyatu.}
\end{quote}

\begin{quote}
A lizard ran across the wall.
\end{quote}

The plural is \ta{பல்லிகள்}.

\begin{cousinweb}
\textbf{\ta{பல்லி}} (\emph{palli}) is native Dravidian (DEDR 3994), \textquotedblleft{}a wall lizard, a gecko\textquotedblright{}.
\end{cousinweb}

\begin{grammarlens}[title={What you've built}]
One more word for getting around.
\end{grammarlens}

\subsection*{Your turn}
\begin{itemize}
\raggedright
  \item \emph{Say it:} \emph{palli}
  \item \emph{Say it:} \emph{palli}, once more
  \item \emph{Say it:} \emph{cuvaril oru palli ōṭiyatu.} — A lizard ran across the wall.
  \item \emph{Recall:} say the Tamil for a koel, then the Tamil for a mongoose, then say \emph{palli} again
\end{itemize}

\begin{tcolorbox}[breakable,colback=teal!4,colframe=teal!35!black,title={Before you move on}]
What does \ta{பல்லி} mean? (\textquotedblleft{}A lizard; a house gecko\textquotedblright{}.) Say it once more.
\end{tcolorbox}
\section[\texorpdfstring{\ta{தேள்} (tēḷ)}{தேள் (tēḷ)}]{\ta{தேள்} (tēḷ) — a scorpion}

\label{lesson:TA-C438-tel}



\begin{quote}
Before the new one: say the Tamil for a mongoose, then the Tamil for a lizard.
\end{quote}

\subsection*{You'll want to know: \ta{தேள்}}
\textbf{\ta{தேள்}} — \emph{tēḷ} — \textquotedblleft{}a scorpion\textquotedblright{}.

\begin{quote}
\ta{பாறையின்} \ta{கீழே} \ta{ஒரு} \ta{தேள்} \ta{இருந்தது}.
\end{quote}

\begin{quote}
\emph{pāṟaiyiṉ kīḻē oru tēḷ iruntatu.}
\end{quote}

\begin{quote}
There was a scorpion under the rock.
\end{quote}

Its sting is very painful.

\begin{cousinweb}
\textbf{\ta{தேள்}} (\emph{tēḷ}) is native Dravidian (DEDR 3470), \textquotedblleft{}a scorpion\textquotedblright{}.
\end{cousinweb}

\begin{grammarlens}[title={What you've built}]
One more word for getting around.
\end{grammarlens}

\subsection*{Your turn}
\begin{itemize}
\raggedright
  \item \emph{Say it:} \emph{tēḷ}
  \item \emph{Say it:} \emph{tēḷ}, once more
  \item \emph{Say it:} \emph{pāṟaiyiṉ kīḻē oru tēḷ iruntatu.} — There was a scorpion under the rock.
  \item \emph{Recall:} say the Tamil for a mongoose, then the Tamil for a lizard, then say \emph{tēḷ} again
\end{itemize}

\begin{tcolorbox}[breakable,colback=teal!4,colframe=teal!35!black,title={Before you move on}]
What does \ta{தேள்} mean? (\textquotedblleft{}A scorpion\textquotedblright{}.) Say it once more.
\end{tcolorbox}
\section[(dialogue)]{Looking back: from Canals and Wells to Herons and Koels}

\label{lesson:TA-R438-review-432-438}



\begin{quote}
Say the Tamil for a body of water and a canal.
\end{quote}

\subsection*{Your turn}
Say these aloud:

\begin{itemize}
\raggedright
  \item a body of water, a canal, a dam, a water well, to slosh
  \item a peak, to trek, a volcano, a glacier, a plain
  \item a meadow, to graze, a grove, a palmyra palm, a bush
  \item the coast, a bay, a cape, a salt pan, a pebble
  \item a whale, a shark, a crab, a seashell, coral
  \item wildlife, a sanctuary, to migrate, to hunt, a species
  \item a heron, a koel, a mongoose, a lizard, a scorpion
\end{itemize}

\begin{tcolorbox}[breakable,colback=teal!4,colframe=teal!35!black,title={Before you move on}]
A body of water? (\textbf{\ta{நீர்நிலை}}, \emph{nīrnilai}.) A canal? (\textbf{\ta{கால்வாய்}}, \emph{kālvāy}.) A dam? (\textbf{\ta{அணைக்கட்டு}}, \emph{aṇaikkaṭṭu}.) A water well? (\textbf{\ta{கிணறு}}, \emph{kiṇaṟu}.) To slosh? (\textbf{\ta{தளும்பு}}, \emph{taḷumpu}.) A peak? (\textbf{\ta{சிகரம்}}, \emph{cikaram}.) To trek? (\textbf{\ta{மலையேறு}}, \emph{malaiyēṟu}.) A volcano? (\textbf{\ta{எரிமலை}}, \emph{erimalai}.) A glacier? (\textbf{\ta{பனிப்பாறை}}, \emph{paṉippāṟai}.) A plain? (\textbf{\ta{சமவெளி}}, \emph{camaveḷi}.) A meadow? (\textbf{\ta{புல்வெளி}}, \emph{pulveḷi}.) To graze? (\textbf{\ta{மேய்}}, \emph{mēy}.) A grove? (\textbf{\ta{தோப்பு}}, \emph{tōppu}.) A palmyra palm? (\textbf{\ta{பனை}}, \emph{paṉai}.) A bush? (\textbf{\ta{புதர்}}, \emph{putar}.) The coast? (\textbf{\ta{கடலோரம்}}, \emph{kaṭalōram}.) A bay? (\textbf{\ta{வளைகுடா}}, \emph{vaḷaikuṭā}.) A cape? (\textbf{\ta{முனை}}, \emph{muṉai}.) A salt pan? (\textbf{\ta{உப்பளம்}}, \emph{uppaḷam}.) A pebble? (\textbf{\ta{கூழாங்கல்}}, \emph{kūḻāṅkal}.) A whale? (\textbf{\ta{திமிங்கலம்}}, \emph{timiṅkalam}.) A shark? (\textbf{\ta{சுறா}}, \emph{cuṟā}.) A crab? (\textbf{\ta{நண்டு}}, \emph{naṇṭu}.) A seashell? (\textbf{\ta{சிப்பி}}, \emph{cippi}.) Coral? (\textbf{\ta{பவளம்}}, \emph{pavaḷam}.) Wildlife? (\textbf{\ta{வனவிலங்கு}}, \emph{vaṉavilaṅku}.) A sanctuary? (\textbf{\ta{சரணாலயம்}}, \emph{caraṇālayam}.) To migrate? (\textbf{\ta{இடம்பெயர்}}, \emph{iṭampeyar}.) To hunt? (\textbf{\ta{வேட்டையாடு}}, \emph{vēṭṭaiyāṭu}.) A species? (\textbf{\ta{இனம்}}, \emph{iṉam}.) A heron? (\textbf{\ta{கொக்கு}}, \emph{kokku}.) A koel? (\textbf{\ta{குயில்}}, \emph{kuyil}.) A mongoose? (\textbf{\ta{கீரி}}, \emph{kīri}.) A lizard? (\textbf{\ta{பல்லி}}, \emph{palli}.) A scorpion? (\textbf{\ta{தேள்}}, \emph{tēḷ}.)
\end{tcolorbox}
